% !TeX program = xelatex
% UTF-8. Compile with XeLaTeX from this directory or the repository root.
% Define \HypernatremiaInNotebook before inputting into a XePersian document.
% Source: Google Drive Notebook/hypernatremia, seven handwritten pages.
% pic01.jpg = 20261016_053657253.jpg (15L57LD2TivUY28o9wLhUnTGuMBoB4kxJ)
% pic02.jpg = 20261016_053702459.jpg (1jR6lu2RxChu4REErEULJyG1I6Hbyg70H)
% pic03.jpg = 20261016_053706088.jpg (1bPh2_yjAe3YzQKPJ69F_XW6QJSb6H0PZ)
% pic04.jpg = 20261016_053712733.jpg (1DTrFDmmRpKzjYIlppiZmtFWy9BEAL39d)
% pic05.jpg = 20261016_053716019.jpg (1u2tL-W5hnD7XKYmkHrA_fFgYsB42WX8u)
% pic06.jpg = 20261016_053721086.jpg (14z2lg9dXDUxkt3m3x4CCL_JGV_URAjKp)
% pic07.jpg = 20261016_053723860.jpg (1JwSYYTsRKmBTljXn3X9SCsF2s8_epSDU)
\ifdefined\HypernatremiaInNotebook
  \let\HNfinish\relax
\else
  \def\HNfinish{\end{document}}
  \documentclass[12pt,a4paper]{article}
  \usepackage[margin=22mm]{geometry}
  \usepackage{amsmath,amssymb}
  \usepackage{graphicx}
  \usepackage{xepersian}
  \IfFontExistsTF{Amiri}{\settextfont{Amiri}}{\settextfont{DejaVu Sans}}
  \setlatintextfont{DejaVu Serif}
  \setlength{\parindent}{0pt}
  \setlength{\parskip}{5pt}
  \setlength{\emergencystretch}{3em}
  \begin{document}
\fi
\providecommand{\HNmath}[1]{\lr{\(#1\)}}
\providecommand{\HNunclear}[1]{\textbf{[خوانش نامطمئن یا متن بریده: #1]}}
\providecommand{\HNpage}[2]{\clearpage\section*{تصویر #1: \lr{#2}}}

\begin{center}
{\Large\bfseries \lr{Hypernatremia in children}}
\end{center}
\textbf{یادداشت رونویسی:} متن به ترتیب هفت تصویر و با حفظ عبارت‌ها، اعداد، اختصارات و روابط فلش‌ها بازنویسی شده است. موارد ناخوانا یا بریده در لبهٔ عکس مشخص شده‌اند؛ متن منبع از نظر پزشکی تصحیح نشده است. سربرگ چاپی همهٔ صفحات «موضوع» (خالی) و «تاریخ / /» (خالی) است؛ شماره‌های دست‌نویس صفحات به ترتیب ۱ تا ۷ هستند. عنوان دست‌نویس بالای صفحهٔ اول «\lr{Hypernatremia in children}» است. اصل هفت تصویر در پیوست همین سند آمده است.

\HNpage{۱}{pic01.jpg}
\begin{itemize}
\item در کودکان: \HNmath{150<\mathrm{Na}}.
\item یک اختلال الکترولیت ناشایع در کودکان است.
\item در موارد \lr{Severe} می‌تواند آسیب نورولوژیک جدی به همراه داشته باشد؛ خصوصاً در موارد \lr{Acute} و \lr{Rapid}.
\end{itemize}
\subsection*{اپیدمیولوژی}
\begin{itemize}
\item ۰٫۰۴\% کودکان بستری انتظار داریم. در آزمایش: ۱٫۴\%. \HNunclear{عبارت «در آزمایش» در سمت چپ سطر}
\item در نوزادان ۰٫۴\% (۱۰ برابر).
\item شایع‌ترین علت دهیدریشن شدید در زمینهٔ \lr{GE} است؛ یا بیماری سیستمیک. \HNunclear{انتهای «سیستمیک» در لبهٔ چپ تصویر}
\end{itemize}
\subsection*{پاتوفیزیولوژی}
\begin{itemize}
\item به علت ناتوان بودن در جایگزینی مایعات از دست رفته (به طور شایع در \lr{GE}) رخ می‌دهد.
\item در این بیماران، آب بیشتر از الکترولیت از دست می‌رود.
\item دقت: در صورتی که کودک نتواند مایع از دست رفته را جایگزین کند هایپرناترمی رخ می‌دهد.
\item در موارد نادری به علت مسمومیت با سدیم رخ می‌دهد.
\end{itemize}
\subsection*{اتیولوژی}
\begin{itemize}
\item از دست دادن مایع که جایگزین نشود.
\item دریافت نمک زیاد و عدم دریافت آب کافی.
\end{itemize}

\HNpage{۲}{pic02.jpg}
\subsection*{از دست دادن آب}
\begin{itemize}
\item \lr{GI}: شایع‌ترین علت هایپرناترمی در کودکان.
\item \lr{AGE} در زمینهٔ روتاویروس.
\item استفراغ و ساکشن \lr{NGT}.
\item ادراری: نقص در تغلیظ ادرار معمولاً در زمینهٔ کمبود یا مقاومت \lr{ADH}.
\end{itemize}
\begin{latin}
\textbf{Arginine Vasopressin Deficiency (AVP-D)}
\end{latin}
\begin{itemize}
\item قبلاً به آن \lr{DI} مرکزی گفته می‌شد.
\item تولید یا آزادسازی ناکافی \lr{ADH}.
\item اتیولوژی‌های مختلف.
\end{itemize}
\begin{latin}
\textbf{Arginine Vasopressin Resistance (AVP-R)}
\end{latin}
\begin{itemize}
\item قبلاً به آن \lr{DI} نفروژنیک می‌گفتند.
\item پاسخ ناکافی توبول کلیه به \lr{ADH}.
\item اتیولوژی‌های مختلف.
\end{itemize}
\begin{latin}
\begin{itemize}
\item Congenital AVP-R (X-linked).
\item Acquired AVP-R: Drug Toxicity \(\longrightarrow\) Lithium.
\item Kidney diseases.
\end{itemize}
\end{latin}
\subsection*{دیورز اسموتیک}
مانیتول و گلوکز (در \lr{DKA}).
\subsection*{\lr{Skin loss}}
غلظت \lr{Na} و \lr{K} عرق کمتر از پلاسما است.

\HNpage{۳}{pic03.jpg}
\subsection*{عدم دریافت کافی مایع}
\begin{itemize}
\item آدیپسی، هایپودیپسی.
\item اختلال در مکانیسم تشنگی: هولوپروزنسفالی، کرانیوفارنژیوما.
\end{itemize}
\subsection*{دریافت زیاد از حد نمک}
\begin{itemize}
\item عموماً ایاتروژنیک است.
\item تجویز سدیم بی‌کربنات یا سالین هایپرتونیک.
\item زمانی که \lr{loss} هایپوتونیک با محلول ایزوتونیک جایگزین شود.
\end{itemize}
\begin{latin}
DKA, Post AKI, Suction NGT.
\end{latin}
یک قاشق چای‌خوری نمک \lr{100 mEq Na} دارد که در یک بچهٔ \lr{10 kg} می‌تواند سدیم را \lr{15 mEq/L} بالا ببرد؛ \lr{OED} ممکن است رخ دهد.
\subsection*{تظاهرات بالینی هایپرناترمی}
\begin{itemize}
\item \lr{Acute}: عموماً علائم نورولوژیک.
\item شدت علائم با شدت هایپرناترمی مرتبط است.
\item تحریک‌پذیری، بی‌قراری، ضعف، استفراغ، توییچینگ عضلانی، تب. \HNunclear{«عضلانی» به صورت حاشیه‌ای در لبهٔ چپ؛ بخشی از حروف بریده است}
\item در نوزادان گریهٔ \lr{high pitched} و تاکی‌پنه.
\item موارد \lr{Severe}، در افزایش حاد به بالای ۱۶۰: \lr{MSC}، لتارژی، تشنج، \lr{SAH}، \lr{OED}، عوارض برگشت‌ناپذیر.
\end{itemize}

\HNpage{۴}{pic04.jpg}
\begin{itemize}
\item \lr{Chronic} (بیش از یک روز): اغلب بی‌علامت.
\item به دلیل انطباق مغز که ۱–۳ روز طول می‌کشد.
\end{itemize}
سطح سدیم خون مویرگی \lr{2--3 mEq/L} کمتر از خون وریدی است؛ در موارد اندازه‌گیری‌شده با دستگاه \lr{ABG} هم کمتر است.
\subsection*{هایپرناترمی کاذب}
\begin{itemize}
\item ورزش شدید، تشنج.
\item سدیم سرمی چند دقیقه \lr{10--15 mEq/L} بالا می‌رود.
\item به علت ورود مایع به داخل سلول.
\item ظرف \lr{5--15 min} به سطح \lr{N/L} برمی‌گردد.
\end{itemize}
در نوزاد \lr{ill} با هایپوآلبومینمی (\HNmath{<30\,\mathrm{g/L}}): به صورت کاذب سدیم \HNunclear{علامت افزایش در حاشیهٔ چپ، بخشی از نوشته بریده است}؛ با \lr{point of care} اندازه‌گیری شود.
\subsection*{ارزیابی کودک با هایپرناترمی}
\begin{itemize}
\item تمرکز بر تعیین اتیولوژی است.
\item اگر بیمار \lr{Severely ill} باشد باید اول مایع‌درمانی کنیم.
\end{itemize}
در اغلب موارد، شرح حال بیمار علت هایپرناترمی را نشان می‌دهد.

\HNpage{۵}{pic05.jpg}
\begin{itemize}
\item شرح حال \lr{GE} و اسهال.
\item شرح حال پلی‌اوری و ادرار رقیق.
\item اختلال نورولوژیک.
\end{itemize}
\subsection*{ارزیابی آزمایشگاهی}
\begin{itemize}
\item قبل از مایع‌درمانی، «مایع» (خط‌خورده) آزمایش‌ها باید گرفته شوند. \textit{(کلمهٔ «مایع» پس از ویرگول در منبع خط خورده است.)}
\item اگر با شرح حال علت مشخص نشد: مقایسهٔ اسمولالیتهٔ ادرار و سرم.
\item \HNmath{U<P}: اختلال تغلیظ ادرار؛ سنترال \lr{DI}، نفروژنیک \lr{DI}، بیماری کلیه و دیورز \HNunclear{ادامهٔ عبارت در لبهٔ چپ تصویر بریده است؛ چند حرف حاشیه‌ای نیز کامل خوانده نمی‌شود}.
\item \HNmath{P<U}: \lr{GI}، تعریق، تب، مسمومیت با نمک.
\end{itemize}
\subsection*{سایر آزمایش‌ها}
\begin{itemize}
\item \lr{BUN}، \lr{Cr}: ارزیابی عملکرد کلیه.
\item محاسبهٔ فراکشنال اکسکریشن سدیم: \lr{FENa}.
\item در \lr{GI} معمولاً هایپرناترمیک هایپوولمی: سدیم ادرار \HNmath{<25}.
\item در مسمومیت با نمک: سدیم ادرار \HNmath{200<\text{urine Na}}.
\end{itemize}
\begin{latin}
\(\mathrm{FENa}>2\%\) in a volume-replete patient is strongly suggestive of hypernatremia caused by \underline{water loss}.
\end{latin}

\HNpage{۶}{pic06.jpg}
\subsection*{درمان هایپرناترمی}
\begin{itemize}
\item جایگزینی مایع از دست رفته.
\item جلوگیری از از دست رفتن مایع بیشتر.
\end{itemize}
در کیس‌هایی که مایع ایزوتونیک هم عقب هستند، دادن \lr{N/S} برای بازگشت گردش خون محیطی مقدم بر اصلاح هایپرناترمی است.
\subsection*{به این سؤالات پاسخ بدهیم}
\begin{itemize}
\item وضعیت حجمی فعلی بیمار چگونه است؟ نیاز اورژانسی به احیا دارد؟
\item میزان دفیسیت مایع چقدر است؟
\item با چه \lr{rate} هایپرناترمی باید اصلاح شود؟
\item آیا به طور همزمان بیمار دارد مایع از دست می‌دهد؟
\item علت زمینه‌ای هایپرناترمی چیست؟
\end{itemize}
\subsection*{بیماری که نیازمند احیا اورژانسی است}
\begin{itemize}
\item پایداری قلبی عروقی باید به سرعت بازگردانده شود.
\item بیمار با هایپوولمی متوسط تا شدید: مایع ایزوتونیک.
\item نباید بیش از حد مایع بدهیم: ادم مغزی.
\end{itemize}
\subsection*{محاسبهٔ دفیسیت مایع}
دقت شرایط اورژانسی نیست؛ میزان \lr{free water} برای اصلاح سدیم را محاسبه می‌کنیم.

\HNpage{۷}{pic07.jpg}
\begin{latin}
\[
\begin{aligned}
&\text{Free water deficit in L}\\
&\quad=\underbrace{0.6\times\text{BW in kg}}_{\mathrm{TBW}}
\times\left(\frac{\mathrm{Na}}{140}-1\right).
\end{aligned}
\]
\[
\begin{aligned}
&\text{Free water deficit in mL}\\
&\quad=(4\,\mathrm{mL/kg})\times\text{BW in kg}\\
&\qquad\times\text{Desired change in Na}.
\end{aligned}
\]
\end{latin}
\HNunclear{انتهای «Desired change in Na» پشت حلقهٔ صحافی و لبهٔ راست عکس پنهان است؛ تا «N» دیده می‌شود.}
\begin{itemize}
\item اگر چه هایپوولمی هم داشت \lr{N/S} می‌دهیم.
\item اگر \lr{H/S} بدهیم برای \lr{500 mL} مایع \lr{free water} باید \lr{1 lit}، \lr{H/S} بدهیم.
\item مایع خوراکی معادل همان حجم \lr{free water} است.
\end{itemize}
\subsection*{نرخ اصلاح هایپرناترمی}
\begin{itemize}
\item در موارد مزمن یا حاد شدید (\HNmath{>160}).
\item حداکثر سرعت اصلاح \lr{0.5 mEq/L/h} (\lr{10--12 mEq/L/day}).
\end{itemize}

\clearpage
\section*{پیوست: اصل تصاویر}
\IfFileExists{images/pic01.jpg}{\def\HNimagebase{images/}}{%
  \def\HNimagebase{pediatrics/ped-nephrology/hypernatremia/images/}}
\newcommand{\HNoriginal}[2]{%
  \clearpage\subsection*{اصل تصویر #1: \lr{pic#2.jpg}}%
  \begin{center}%
  \includegraphics[width=\linewidth,height=0.87\textheight,keepaspectratio]{\HNimagebase pic#2.jpg}%
  \end{center}}
\HNoriginal{۱}{01}
\HNoriginal{۲}{02}
\HNoriginal{۳}{03}
\HNoriginal{۴}{04}
\HNoriginal{۵}{05}
\HNoriginal{۶}{06}
\HNoriginal{۷}{07}
\HNfinish
